% p352_2; source PDF page 352; manually verified editable table.
\begin{tikzpicture}[x=1cm,y=-1cm,inner sep=0pt,outer sep=0pt]
\path[use as bounding box] (0,0) rectangle (16.000,12.610);
\draw[line width=.4pt] (0.000,0.000) rectangle (2.800,0.865);
\node[align=center,text width=2.540cm,font={\linespread{1}\fontsize{10}{13.50}\selectfont\bfseries}] at (1.400,0.432) {罪名};
\draw[line width=.4pt] (2.800,0.000) rectangle (9.750,0.865);
\node[align=center,text width=6.690cm,font={\linespread{1}\fontsize{10}{13.50}\selectfont\bfseries}] at (6.275,0.432) {构成要件};
\draw[line width=.4pt] (9.750,0.000) rectangle (16.000,0.865);
\node[align=center,text width=5.990cm,font={\linespread{1}\fontsize{10}{13.50}\selectfont\bfseries}] at (12.875,0.432) {要点};
\draw[line width=.4pt] (0.000,0.865) rectangle (2.800,8.080);
\node[align=center,text width=2.540cm,font={\linespread{1}\fontsize{10}{13.50}\selectfont}] at (1.400,4.473) {非法吸收公众存款罪};
\draw[line width=.4pt] (2.800,0.865) rectangle (9.750,8.080);
\node[align=left,text width=6.690cm,font={\linespread{1}\fontsize{10}{13.50}\selectfont}] at (6.275,4.473) {1.非法性。这是指未经批准。\\{}2.公开性。这是指通过媒体等途径向社会公开宣传。\\{}3.利诱性。这是指承诺在一定期限内还本付息或者给付回报。\\{}4.公众性。这是指向社会公众即社会不特定对象吸收资金。\\{}5.存款性。这是指，行为人吸收存款是想用于放贷等金融活动，由此吸收的资金才能称为吸收存款。如果吸收资金用于正当的生产经营活动，则仅属于吸收资金，不属于吸收存款};
\draw[line width=.4pt] (9.750,0.865) rectangle (16.000,8.080);
\node[align=left,text width=5.990cm,font={\linespread{1}\fontsize{10}{13.50}\selectfont}] at (12.875,4.473) {1.向亲友或者单位内部人员吸收资金的，不属于向社会公众吸收存款；但如果明知亲友或者单位内部人员又向不特定对象吸收资金而予以放任的，属于向社会公众吸收存款。（2014年第14题）\\{}2.通过传销手段实施，同时触犯组织、领导传销活动罪和非法吸收公众存款罪或集资诈骗罪的，属于牵连犯，择一重罪论处};
\draw[line width=.4pt] (0.000,8.080) rectangle (2.800,10.595);
\node[align=center,text width=2.540cm,font={\linespread{1}\fontsize{10}{13.50}\selectfont}] at (1.400,9.338) {集资诈骗罪};
\draw[line width=.4pt] (2.800,8.080) rectangle (9.750,10.595);
\node[align=left,text width=6.690cm,font={\linespread{1}\fontsize{10}{13.50}\selectfont}] at (6.275,9.338) {通过非法集资的方式实施的诈骗罪};
\draw[line width=.4pt] (9.750,8.080) rectangle (16.000,10.595);
\node[align=left,text width=5.990cm,font={\linespread{1}\fontsize{10}{13.50}\selectfont}] at (12.875,9.338) {1.集资诈骗罪有非法占有目的。非法吸收公众存款罪没有非法占有目的。\\{}2.集资后不用于生产经营，或携带集资款逃匿，视为具有非法占有目的};
\draw[line width=.4pt] (0.000,10.595) rectangle (2.800,12.610);
\node[align=center,text width=2.540cm,font={\linespread{1}\fontsize{10}{13.50}\selectfont}] at (1.400,11.603) {擅自发行股票、公司、企业债券罪};
\draw[line width=.4pt] (2.800,10.595) rectangle (9.750,12.610);
\node[align=left,text width=6.690cm,font={\linespread{1}\fontsize{10}{13.50}\selectfont}] at (6.275,11.603) {未经批准，向社会不特定对象发行、以转让股权等方式变相发行股票或者公司、企业债券};
\draw[line width=.4pt] (9.750,10.595) rectangle (16.000,12.610);
\node[align=left,text width=5.990cm,font={\linespread{1}\fontsize{10}{13.50}\selectfont}] at (12.875,11.603) {针对特定对象发行，特定对象人数超过200人的，也构成本罪};
\end{tikzpicture}
